\title{Byte Latent Transformer: Patches Scale Better Than Tokens}

\begin{abstract}
We present the Byte Latent Transformer ({\textsc BLT}), a novel large language model architecture operating at the byte level that, for the first time, achieves parity with token-based LLMs in terms of performance when scaled, while offering marked gains in inference efficiency and robustness.
{\textsc BLT} transforms raw bytes into adaptively sized patches, which function as the main computational primitives. The model divides these patches according to the entropy of the upcoming byte, dedicating increased computation and model resources to regions of higher data complexity. We conduct the initial flop-controlled scaling analysis for byte-level language models, evaluating models up to 8B parameters across 4T training bytes. These experiments verify that scaling LLMs directly on byte sequences—without relying on a predetermined vocabulary—is feasible. Dynamic allocation of longer patches over predictable segments not only enhances training and inference efficiency, but also yields qualitative improvements in reasoning abilities and generalization to rare cases. In summary, at a fixed inference budget, {\textsc BLT} achieves stronger scaling behavior than tokenization-based approaches by jointly expanding both model and patch sizes.
\end{abstract}

\section{Introduction}
\label{section:intro}

We present the Byte Latent Transformer~(\textbf{BLT}), a novel architecture that eliminates the need for explicit tokenization by learning directly from unprocessed byte sequences. For the first time, this approach achieves competitive results with established tokenization-based systems at large scale, while offering considerable gains in computational efficiency and robustness (\S\ref{sec:robustness}). In the prevailing paradigm, large language models (llms) are almost fully end-to-end, yet rely on tokenization—a pre-processing technique that heuristically groups bytes into a fixed vocabulary of tokens. This step imposes a representational bias, affecting compression strategies and resulting in limitations such as sensitivity to domain and modality~\citep{dagangetting}, heightened vulnerability to input perturbations~(\S\ref{sec:robustness}), insufficient recognition of orthographic features~\citep{edman2024cute}, and disparities across languages~\citep{liang2023xlm,petrov2024language,limisiewicz2024myte}.

In earlier research, tokenization played a crucial role, as training large language models (llms) directly on byte-level data incurs significant computational expenses caused by the increased sequence lengths~\citep{xue2022byt5}. Prior approaches have addressed this issue through the adoption of more resource-efficient self-attention mechanisms~\citep{el2020characterbert, clark2022canine} or by exploring architectures that eliminate attention entirely~\citep{wang2024mambabyte}~(\S\ref{section:related_work}). Nonetheless, these solutions mainly benefit the training of \textit{smaller models}. When scaling up, the primary computational bottleneck in Transformers is typically the large feed-forward network layers that process every byte, rather than the cost associated with the attention mechanism.

In order to optimize compute distribution, we introduce a dynamic, trainable approach for assembling bytes into \textit{patches}~(\S\ref{section:patching}), along with a novel architecture capable of integrating both byte-level and patch-level data. In contrast to tokenization, {\textsc BLT} does not depend on a predefined set of patches. Instead, any subset of bytes can be transformed into latent patch representations through efficient, trainable encoder and decoder components. Our findings demonstrate that this approach allows for a \textit{greater} efficiency in compute utilization compared to models that rely on tokenization.

Conventional tokenization-based LLMs uniformly distribute computational resources across all tokens, prioritizing consistency at the expense of optimal efficiency. This is due to token formation relying on compression strategies that do not necessarily align with the true complexity involved in making predictions. Our proposed approach is grounded in the principle of adaptively allocating computational effort where it is most warranted. For instance, predicting the conclusion of most words does not require the capacity of a full-scale transformer, as such tasks typically present less complexity and lower entropy than generating the initial word of a new sentence. This adaptive compute distribution is exemplified in the {\textsc BLT} design~(\S\ref{section:architecture}), which comprises three transformer modules: a pair of lightweight, byte-level \textit{local models}, complemented by a sizable, global \textit{latent transformer}~(\autoref{fig:arch_high_level}). To achieve effective dynamic compute allocation, {\textsc BLT} organizes bytes into patches by analyzing the entropy in next-byte predictions. This process yields contextual patches containing bytes that share a similar level of informational content.

We introduce the inaugural study of scaling byte-level models constrained by training flops, investigating models with sizes reaching 8B parameters and trained on datasets containing up to 4T bytes. Our results demonstrate the feasibility of training large-scale end-to-end byte-based models, dispensing entirely with the need for fixed-vocabulary tokenization. Notably, {\textsc BLT} achieves parity with Llama 3 in terms of flop-normalized training performance,\footnote{We assess model compute expenditure via the total count of Floating Point OPerations (flops) required.} while also reducing inference-time flops by as much as 50\% (\S\ref{section:scaling}). 

Moreover, we find that operating at the raw byte level confers meaningful advantages in handling the statistical long-tail within data distributions. {\textsc BLT} outperforms tokenizer-dependent approaches in robustness to noisy input and demonstrates superior comprehension at the character level—exhibited through improved orthographic recognition, phonological discrimination, and performance on tasks such as low-resource machine translation (\S\ref{sec:robustness}).

Additionally, {\textsc BLT} permits the concurrent enlargement of both model and patch sizes while adhering to a constant inference flop constraint. Utilization of longer patch sizes generally leads to compute savings that may instead be allocated towards expanding the global latent transformer's capacity, as its invocation frequency is subsequently reduced. Through experiments controlling for inference flops (\autoref{fig:fixed_inference_scaling}), we identify more favorable scaling behavior compared to traditional tokenization-based models.

To conclude, our work contributes the following: 1) We present {\textsc BLT}, a byte latent LLM model designed to allocate computational resources dynamically, thereby enhancing flop efficiency; 2) We demonstrate that our approach attains parity with Llama 3 at scales up to 8B parameters, in terms of training flop consumption, while also allowing users to exchange slight evaluation performance reductions for up to 50\% gains in flop efficiency; 3) {\textsc BLT} architectures enable a novel scaling strategy for LLMs, where the model size may be increased without exceeding a predetermined inference cost; and 4) We show that {\textsc BLT} models exhibit heightened robustness to input perturbations and better sensitivity to the sub-word information embedded in byte sequences—a feature that token-based models typically overlook. The training and inference code for {\textsc BLT} is publicly available at~\url{https://github.com/facebookresearch/blt}.

\section{Patching: From Individual Bytes to Groups of Bytes}
\label{section:patching}

Dividing byte sequences into \textit{patches} enables {\textsc BLT} to flexibly distribute computational resources based on the surrounding context. Several approaches for segmenting bytes into patches are illustrated in Figure~\ref{fig:patching_types}. More formally, the patching function $f_p$ partitions a byte sequence $\pmb{x} = \{x_i \mid i=1, \ldots, n\}$ of length $n$ into a series of $m < n$ patches $\pmb{p} = \{p_j \mid j=1, \ldots, m\}$, assigning each $x_i$ a value in \{0,1\}, where a 1 marks the beginning of a new patch. For both models built on tokens and those using patches, the main contributor to computational complexity is the number of steps required by the Transformer core. Within {\textsc BLT}, this corresponds to the count of patches necessary to represent the input data given a specific patching approach. As such, the average patch length—simply referred to as \textit{patch size}—serves as the principal metric influencing computational expense in both training and inference phases for a chosen patching method~(\S\ref{section:flops}).

Subsequently, we introduce three kinds of patching functions: static patching based on a fixed byte count per patch~(\S\ref{section:static-patch}), patching at whitespace boundaries~(\S\ref{section:white-patch}), and adaptive patching that leverages entropy measures from a lightweight byte-level language model~(\S\ref{section:dyn-patch}). We conclude with a discussion on incremental patching methods and examine distinctions between tokenization and patching~(\S\ref{section:bpe-patch}).

\subsection{Strided Patching Every K Bytes}
\label{section:static-patch}

A common approach for aggregating bytes involves partitioning them into uniform patches of size $k$, as implemented in MegaByte~\citep{yu2023megabyte}. This method benefits from a fixed stride, which simplifies both model training and inference, while also providing an intuitive lever for adjusting the average patch size, thereby facilitating precise control over computational resources (flop cost). Despite these advantages, using a fixed patching scheme presents notable limitations. Primarily, the allocation of computational effort is not adaptive: computational steps may be inefficiently spent on less informative regions (e.g., whitespace in code), or conversely, insufficient attention may be given to segments where byte content is information-rich, such as in mathematical expressions. Additionally, this rigid approach can result in irregular and context-insensitive segmentation of similar byte strings, with identical words potentially divided inconsistently.

\subsection{Space Patching}
\label{section:white-patch}

\citet{slagle2024spacebyte} introduces a straightforward yet potent modification to strided patching, wherein new patches are initiated immediately after any space-like bytes\footnote{
    Space-like bytes are defined as any byte that is not a latin character, digit, or utf-8 continuation byte. In addition, each patch must contain at least one non space-like byte. }, which often serve as natural demarcations for linguistic segments in a variety of languages. Under Space patching, an additional transformer step (i.e., increased computational effort measured in flops) is dedicated to model each word individually. This design guarantees consistent treatment of words within and across sequences, while also ensuring computational resources are expended on making challenging predictions that frequently arise after space-like bytes. Consider the example: generating the initial byte of the answer to ``Who composed the Magic Flute?\uline{\hspace{.6em}}'' is substantially more complex than producing subsequent bytes, as the initial character of the answer dramatically narrows down the possible continuations—rendering the rest of the answer, such as ``Mozart'', much simpler to anticipate. Nonetheless, space patching has limitations: it does not universally adapt to all languages or domains and, crucially, lacks the flexibility to adjust patch sizes. In the following, we present an alternative patching technique inspired by the observation that the first bytes of words generally pose the greatest predictive challenge, and one that naturally enables control over patch sizing.

\subsection{Entropy Patching: Using Next-Byte Entropies from a Small Byte LM}
\label{section:dyn-patch}

Instead of utilizing a heuristic based on specific rules like whitespace, our method adopts a data-centric strategy to detect points of high uncertainty in predicting the next byte. We propose \textit{entropy patching}, a technique that leverages entropy measurements to determine the locations of patch boundaries.

To this end, we fit a compact auto-regressive language model operating at the byte level to the {\textsc BLT} training corpus, and calculate the entropy of subsequent byte predictions according to the language model's probability distribution $p_{e}$ over the byte-level vocabulary $\mathcal{V}$:

\begin{equation}
H(x_i) = \sum_{v \in \mathcal{V}} p_{e}(x_i=v|\pmb{x}_{<i}) \log p_{e}(x_i=v|\pmb{x}_{<i})
\end{equation}

We explore two approaches for detecting patch boundaries based on entropies $H(x_i)$. The first method selects points where the entropy surpasses a fixed global threshold, as depicted in~\autoref{fig:patching}. The second method locates points exhibiting high entropy compared to their immediate predecessor. Alternatively, this latter strategy can be seen as pinpointing positions that disrupt an otherwise roughly monotonically decreasing entropy trend within a patch.

\begin{align*}
    \text{Global Constraint}\hspace{1cm}H(x_t)& > \theta_g\\
    \text{Approx. Monotonic Constraint}\hspace{1cm}H(x_t)& - H(x_{t-1}) > \theta_r
\end{align*}

Patch boundaries are determined in a minimal preprocessing phase that occurs at the data loading stage. This contrasts with the approach of~\citet{nawrot-etal-2023-efficient}, who utilize a classifier specifically trained to infer patch boundaries based on entropy measures. In the course of our experiments~(\S\ref{section:experiments}), we evaluate and compare these two strategies for differentiating between bytes of low and high entropy.

\subsection{The Byte-Pair Encoding (BPE) Tokenizer and Incremental Patching}
\label{section:incremental-patching}
\label{section:bpe-patch}

A wide range of contemporary llms, such as the Llama 3 baseline, employ subword tokenizers like bpe~\citep{gage1994new, sennrich-etal-2016-neural}. In this context, we define ``tokens'' as sequences of bytes sampled from a \textit{finite} vocabulary established ahead of training, distinguishing them from ``patches,'' which are formed from adaptively segmented sequences lacking a predetermined vocabulary. Importantly, unlike patches, tokens do not enable the model to directly utilize the original byte-level information.

A significant advancement introduced by {\textsc BLT} compared to tokenization-based architectures is its reconfiguration of the relationship between vocabulary size and computational requirements. Traditional language models must balance the growth of vocabulary: while enlarging the vocabulary leads to longer tokens and, consequently, reduces the total number of processing steps, it simultaneously expands the output dimensionality of the final projection layer. This inherent compromise restricts tokenization-based methods, making it challenging to realize substantial changes in both token length and the associated inference overhead. As an illustration, Llama 3 opts to raise the average token length from 3.7 to 4.4 bytes; however, this improvement necessitates a fourfold expansion of its embedding table relative to Llama 2.

During generation, {\textsc BLT} must determine at each point in the byte sequence whether it is positioned at a patch boundary, which impacts whether additional computation through the Latent Transformer is necessary. This determination must be made without reliance on bytes that have not yet been generated; that is, the segmentation strategy for patching must operate without knowledge of future bytes. More precisely, a patching scheme $f_p$ adheres to the principle of incremental patching if:
$$f_p(\pmb{x}_{<i}) = f_p(\pmb{x})_{<i}$$
bpe does not constitute an incremental patching method since a given prefix might receive different tokenizations based on the subsequent sequence, violating the above criterion\footnote{Introducing a unique delimiter token at patch boundaries allows bpe to be adapted for incremental patching, though this leads to a longer byte sequence.}.

\section{{\textsc BLT} Architecture}
\label{section:architecture}
The {\textsc BLT} framework incorporates a large-scale, globally-autoregressive language model that processes patch-level representations. Accompanying this central model are two compact, localized components: one to transform byte sequences into patch representations and another to reconstruct bytes from patches, as outlined in~(\autoref{fig:arch_high_level}).

\subsection{Latent Global Transformer Model}

\textit{The Latent Global Transformer} refers to an autoregressive transformer architecture denoted as $\mathcal{G}$, characterized by $l_{\mathcal{G}}$ layers. This model takes as input a sequence of latent patch representations, $p_j$, and transforms them into a corresponding sequence of output patch representations, $o_j$. In all discussions, patches are indexed by the subscript $j$, while bytes are indicated by $i$. The architecture utilizes a block-causal attention mask~\citep{dubey2024llama} that confines each token’s attention to the patches up to and including the current one within a document. As the most computationally intensive component in both pre-training and inference stages, the global model's invocation points are crucial; strategically selecting when to apply it enables dynamic adjustment of computational resources depending on the complexity of different regions within the input and output.

\subsection{Local Encoder}
\label{section:local}

\textit{The Local Encoder Model}, symbolized as $\mathcal{E}$, employs a streamlined transformer framework comprising $l_{\mathcal{E}} << l_{\mathcal{G}}$ layers. Its primary objective is the rapid transformation of a sequence of input bytes $b_i$ into rich patch-level representations $p_j$. Unlike conventional transformer architectures, this model incorporates a cross-attention mechanism after every transformer layer, allowing for aggregation of byte-level embeddings into patch representations~(\autoref{fig:crossattn}).

Initially, the bytes $b_i$ are mapped to embedding vectors, $x_i$, via a matrix of shape $\mathbb{R}^{256 \times h_{\mathcal{E}}}$. Optionally, these initial embeddings can be further enriched using hash-embeddings~(\S\ref{sec:hash-emb}). 

A stack of transformer and cross-attention layers arranged alternately~(\S\ref{sec:enc-cross-attn}) subsequently refines these embeddings into the patch representations $p_i$, which serve as input to the global transformer module, $\mathcal{G}$. The attention mask utilized within each transformer block is a \textit{local block causal} mask: each byte attends only to a bounded context of $w_{\mathcal{E}}$ prior bytes. This contextual window may traverse across dynamically-determined patch boundaries but is restricted from extending past the boundaries of the document. Subsequent sections provide in-depth discussion of both the embedding mechanism and the implementation of the cross-attention block.

\subsubsection{Encoder Hash n-gram Embeddings}
\label{sec:ngram-emb}
\label{sec:hash-emb}
An essential aspect of generating strong and informative representations at every position $i$ involves capturing context from previous bytes. Within {\textsc BLT}, this is accomplished by treating the current byte $b_i$ both on its own and as an element within a byte n-gram sequence. Specifically, at each step $i$, we begin by forming byte-grams

\begin{align}
    g_{i,n}=\{b_{i-n + 1},\ldots, b_i\}
\end{align}

for each byte position $i$ and $n$ from three to eight.\footnote{
We omit byte-grams of size $n$ or more when $i<n$. }

Next, we present hash $n$-gram embeddings, wherein each byte $n$-gram is assigned an index in a fixed-size embedding table $E_{n}^{hash}$ using a hash function, for every $n \in\{3,4,5,6,7,8\}$~\citep{bai2010learning}. The obtained $n$-gram embedding is combined with the byte’s own embedding before normalization, and this sum is subsequently fed into the local encoder. We obtain the enhanced embedding

\begin{align}
e_i &= x_i + \sum_{n = 3,...,8}E_{n}^{hash}(\text{Hash}(g_{i,n})) \\
\text{where, Hash}(g_{i,n}) &= \text{RollPolyHash}(g_{i,n}) \% |E_{n}^{hash}|
\end{align}

Each $e_i$ is scaled by dividing by the total number of considered $n$-gram sizes plus one, employing RollPolyHash as specified in Appendix~\ref{appendix:rollpolyhash}. Section~\ref{section:ablations} presents an ablation study examining how varying $n$-gram values and embedding table sizes influence flop-controlled scaling law behaviors. Beyond hash-based $n$-gram embeddings, we also explored frequency-based $n$-gram embedding methods, with comprehensive information available in Appendix~\ref{appendix:ngrams}.

\subsubsection{Encoder Multi-Headed Cross-Attention}
\label{sec:enc-cross-attn}
Our approach draws substantially from the input cross-attention module employed in the Perceiver architecture~\citep{jaegle2021perceiver}. The primary distinction, however, lies in our use of dynamic patch representations serving as the latent entities, replacing the static latent set utilized by the original framework~(\autoref{fig:crossattn}); each patch only engages with the bytes from its own region. The procedure initiates by creating a query vector for each patch $p_j$, achieved through aggregating the corresponding byte-level representations of $p_j$—this aggregation is accomplished by a pooling operation—followed by passing this pooled output through a linear transformation, $\mathcal{E}_{C} \in \mathbb{R}^{h_{\mathcal{E}} \times (h_{\mathcal{E}} \times U_{\mathcal{E}})}$, where $U_{\mathcal{E}}$ stands for the number of heads in the encoder cross-attention layer. Explicitly, defining $f_{\text{bytes}}(p_j)$ as the ordered bytes constituting patch $p_j$, we compute

\begin{align}
P_{0,j} &= \mathcal{E}_{C}(f_\text{bytes}((p_j)), f~ \text{is a pooling function} \\
P_l &= P_{l-1} + W_o\left(\text{softmax}\left(\frac{QK^T}{\sqrt{d_k}}\right)V\right) \\
\text{where } Q_j &= W_q(P_{l-1,j}), K_i = W_k(h_{l-1,i}), V_i = W_v(h_{l-1,i}) \\
h_l &= \text{Encoder-Transformer-Layer}_l(h_{l-1})
\end{align}

Here, $P \in \mathbb{R}^{n_p \times h_{\mathcal{G}}}$ denotes the representations for $n_p$ patches that are input to the global model. These representations are initially formed by aggregating the byte-level embeddings $e_i$ associated with each patch $p_j$. The matrices $W_q$, $W_k$, $W_v$, and $W_o$ are parameter matrices used to project inputs into queries, keys, values, and outputs, respectively; the keys and values are derived by projecting the byte-level representations $h_i$ from the preceding layer (where $e_i$ serves as input in the initial layer). Our attention mechanism incorporates a patch-specific masking approach such that each query vector $Q_j$ can only attend to those key and value vectors that correspond to bytes belonging to patch $j$. Since multi-headed attention is employed for $Q, K$, and $V$, and the patch embeddings typically have higher dimensionality ($h_{\mathcal{G}}$) than $h_{\mathcal{E}}$, we represent $P_l$ as a collection of heads, each with dimension $h_{\mathcal{E}}$, in the cross-attention procedure; these are then concatenated to recover the full $h_{\mathcal{G}}$ dimension. Furthermore, we apply pre-LayerNorm to the queries, keys, and values. No positional encoding is introduced in this cross-attention layer. A residual connection also wraps the cross-attention module to facilitate gradient flow.

\subsection{Local Decoder} 

The local decoder $\mathcal{D}$, like the local encoder, utilizes a transformer architecture designed to be lightweight, comprising $l_{\mathcal{D}}$ layers with $l_{\mathcal{D}} << l_{\mathcal{G}}$. Its purpose is to translate a sequence of global patch embeddings $o_j$ into the output raw bytes $y_i$. This decoder generates each byte in the sequence conditioned on the bytes decoded thus far by leveraging the hidden states computed by the local encoder over the input byte sequence. Processing occurs over $l_{\mathcal{D}}$ successive layers that alternate between cross-attention mechanisms and standard transformer layers. Specifically, the cross-attention mechanism precedes the transformer layer at each stage, enabling the construction of byte-level representations from the higher-level patch embeddings, with subsequent transformer layers operating on this derived byte-level sequence.

\subsubsection{Decoder Multi-headed Cross-Attention} 

Within the decoder's cross-attention mechanism, the assignments for queries and key/values are reversed: the queries consist of the byte-representations, while the keys and values correspond to the patch representations. To initiate cross-attention, the byte-representations are set using the byte embeddings from the terminal encoder layer, namely $h_{l_{\mathcal{E}}}$. For each following decoder layer $l$, the byte-representations $d_{l,i}$ are updated as described below:

\begin{align}
D_0 &= h_{l_{\mathcal{E}}} \\
B_l &= D_{l-1} + W_o\left(\text{softmax}\left(\frac{QK^T}{\sqrt{d_k}}\right)V\right), \\
  &\text{where } Q_i = W_q(d_{l-1,i}),\ K_i = W_k(\mathcal{D}_C(o_j)),\ V_i = W_v(\mathcal{D}_C(o_j)) \\
  D_l &= \text{Decoder-Transformer-layer}_l(B_l)
\end{align}

In this context, the key and value projection matrices $W_k, W_v$ are again used to perform linear transformations and a split operation $\mathcal{D}_C$ on the global model’s final patch representations $o_j$. The query projection matrix $W_q$ is applied to the byte-level representations $d_{l-1}$ from the preceding layer of the decoder transformer (or to $h_{l_{\mathcal{E}}}$ if considering the initial layer). The resulting output is further processed with the projection matrix $W_o$, ensuring that $B$ has dimensions $\mathbb{R}^{h_{\mathcal{D}} \times n_b}$, where $n_b$ denotes the total number of output bytes. Subsequently, the following decoder representations $D_l$ are generated by feeding the output of the cross-attention mechanism, $B$, through a decoder transformer layer. Following the approach used in local encoder cross-attention, the model incorporates multiple attention heads, employs pre-layer normalization, omits positional encodings, and implements a residual connection surrounding the cross-attention component.

\section{Experimental Setup}
\label{section:experiments}
Our experiments are carefully structured to provide a fair comparison between {\textsc BLT} and models based on tokenization, ensuring that {\textsc BLT} does not benefit from access to comparatively longer sequence lengths.

\subsection{Pre-training Datasets} In all model scale configurations explored in this study, pre-training is conducted using two primary datasets: 1) The Llama 2 dataset~\citep{touvron2023llama}, which consists of 2 trillion tokens sourced from a wide range of publicly available data and subsequently subjected to cleaning and filtering processes to ensure higher quality; and 2) {\textsc BLT}-1T: This is a newly constructed dataset containing 1 trillion tokens, which draws from a diversity of public sources and additionally incorporates a subset of the pre-training data made available by Datacomp-LM~\citep{li2024datacomp}. The former serves in scaling law studies to evaluate optimal token counts, following the guidelines set forth by~\cite{dubey2024llama}, in order to inform the architectural choices of {\textsc BLT}. The latter, {\textsc BLT}-1T, is utilized in a comprehensive pre-training run for direct comparison with Llama 3 on a suite of downstream tasks. It should be noted that neither dataset includes information gathered from Meta-related products or services. Additionally, for baseline evaluations involving tokenizer-based models, we implement the Llama 3 tokenizer, which has a 128K token vocabulary, as this configuration achieved better baseline results compared to the Llama 2 tokenizer in our assessments.

\subsection{Entropy Model} For our experiments, the entropy model is constructed as a byte-level language model and is trained using the same data distribution as the overall {\textsc BLT} framework. Unless otherwise noted, our default configuration utilizes a transformer architecture consisting of 100M parameters, 14 transformer layers, a hidden size of 512, and sliding window attention spanning 512 bytes. All other hyperparameter choices mirror those used for both the local and global transformer variants. As further elaborated in \autoref{section:ablations}, we explored a range of model sizes, receptive field lengths, and architectural designs. Notably, when the receptive field is sufficiently restricted, it becomes feasible to represent the learned entropy model as a compact and computationally efficient lookup table.

\subsection{Entropy Threshold and Equalizing Context Length} For models that employ entropy-based patching, we determine a patching threshold that yields a specified average \textit{patch size} over the pretraining data mixture. In contrast to conventional tokenization, {\textsc BLT} permits the \textit{patch size} to be selected with flexibility, substantially affecting the effective context span utilized by the model. To ensure consistent average context lengths and to prevent models with larger patch sizes from gaining an undue advantage, we fix the expected total number of bytes per batch. Consequently, for models with larger patch sizes, we proportionally decrease the number of sequence steps. Specifically, on Llama 2 data, a context length of 8k bytes is adopted, whereas for the {\textsc BLT}-1T dataset, the average context is increased to 16k bytes, with the batch size held constant at 16M bytes on average in both cases.

Although the mean batch size remains unchanged, dynamic patching approaches result in varying byte-to-patch ratios during data batching. In our implementation, the {\textsc BLT} training process arranges batches to maintain a fixed number of patches, thereby circumventing padding within the costly latent transformer stage. To prevent memory usage from spiking due to exceptionally large patches, we pad or, if necessary, truncate byte sequences to 12k bytes for Llama 2 and 24k bytes for {\textsc BLT}-1T datasets during training.

\subsection{Entropy Model Context}
\label{section:entropy-context}
Our experimental observations reveal that entropy patching tends to create increasingly large patches in scenarios involving structured tasks, such as multiple choice formats (as illustrated in the patching procedure on an MMLU example in \autoref{fig:mmlu_patching}), which are typically quite redundant. These differences stem from the lower entropy in the repeated elements within the entropy model context. Accordingly, during the large-scale application of {\textsc BLT}-Entropy with a patch size of 4.5, we refresh the entropy context by introducing new lines and employ the approximate monotonicity constraint, given that it is less susceptible to "entropy drift" arising from fluctuations in context length. Notably, this adjustment influences solely the entropy calculation process; our method for selecting the entropy threshold remains unchanged.

\subsection{FLOPs Estimation} 
\label{section:flops}

Our methodology closely aligns with the transformer flop calculations detailed in Chinchilla~\citep{hoffmann2022training}, incorporating contributions from the feed-forward modules, the qkvo projection operations in self-attention, as well as the computation associated with attention mechanisms and the output projection step. One key modification in our approach is the presumption that the input embedding layer utilizes an efficient lookup table rather than dense matrix multiplication, effectively rendering it a 0-flop process. Consistent with prior studies, we approximate that the backward propagation phase involves double the number of flops as compared to the forward propagation phase.

To determine the flops \textit{per byte} for the {\textsc BLT} architecture, we aggregate the computational costs incurred by the local encoder transformer, global latent transformer, and local decoder transformer. Additionally, the contributions from the cross-attention modules in both the encoder and decoder are included:

\begin{align}
    \text{FL}_{\text{{\textsc BLT}}} &= \frac{\text{Transf. FL}(h_{\mathcal{G}}, l_{\mathcal{G}}, m=n_{ctx}/n_p, V=0)}{n_p}\\
    &\quad +\, \text{Transf. FL}(h_{\mathcal{E}}, l_{\mathcal{E}}, m=w_{\mathcal{E}}, V=0) \\
    &\quad +\, \text{Transf. FL}(h_{\mathcal{D}}, l_{\mathcal{D}}, m=w_{\mathcal{D}}, V=256) \\ 
    &\quad +\, \text{Cross Attn. FL}(h_{\mathcal{E}}, l_{\mathcal{E}}, m=n_p, r=n_p/k) \cdot k/n_p \\ 
    &\quad +\, \text{Cross Attn. FL}(h_{\mathcal{D}}, l_{\mathcal{D}}, m=k, r=k/n_p)
\end{align}

Here, $n_{ctx}$ denotes the input sequence length measured in bytes, and $n_p$ specifies the size of each patch. The parameter $r$ reflects the proportion between the number of queries and the number of keys or values. The factor $k$ represents the ratio between the patch dimension and the byte dimension; it also indicates how many local model components are concatenated to generate the global model representation (as shown with $k=2$ in \autoref{fig:crossattn}). The variable $V$ stands for the vocabulary size utilized during output projection, relevant solely to the local decoder. The attention's context length, represented by $m$, varies based on whether a module is applied to the byte-level or patch-level sequence. We adjust the attention-related flop counts accordingly for each respective model component. Appendix~\ref{section:flops-appendix} presents the specific mathematical expressions for calculating Transformer-FLOPs and Cross-Attention FLOPs.

\subsection{Bits-Per-Byte Estimation} Perplexity is meaningful only when using a consistent tokenizer, as it quantifies the uncertainty associated with individual tokens. To facilitate comparison between models operating at the byte and token levels, we adopt Bits-Per-Byte (BPB), as employed in previous studies~\citep{xue2022byt5, yu2023megabyte, wang2024mambabyte}. BPB serves as a tokenizer-agnostic alternative to perplexity. Concretely:

\begin{align}
    \text{BPB}(x)&= \frac{\mathcal{L}_{CE}(\pmb{x})}{\ln(2)\cdot n_{\text{bytes}}}
\end{align}

where the uncertainty over the data $\pmb{x}$ as measured by the sum of the cross-entropy loss is normalized by the total number of bytes in $\pmb{x}$ and a constant. 

\subsection{Transformer Architecture Hyperparameters} 

For every transformer block within {\textsc BLT}, encompassing both the local and global modules, our configuration closely mirrors that of Llama~3~\citep{dubey2024llama}. Specifically, the feed-forward components utilize the SwiGLU activation function~\citep{swiglu}, while rotary positional embeddings (RoPE)~\citep{RoPe} are applied with a parameter of $\theta=500000$~\citep{xiong2024effective}, limited to self-attention mechanisms. Additionally, RMSNorm~\citep{RMSNorm} is adopted for all normalization layers. Regarding self-attention operations with standard, unchanging attention masks (including \textit{block causal} and \textit{fixed-window block causal}), we implement Flash attention~\citep{dao2022flashattention}. For masks that maintain a fixed width, a window size of 512 is designated. In the case of cross-attention modules, due to the variable and patch-dependent nature of their attention masks, we leverage Flex Attention\footnote{\url{https://pytorch.org/blog/flexattention}} in order to attain fused computation and accelerate the training process considerably.

\subsection{{\textsc BLT}-Specific Hyperparameters} To evaluate the performance of {\textsc BLT} models, we perform experiments focusing on two main aspects: scaling behaviors and downstream task assessments, considering models at varying parameter sizes, specifically 400M, 1B, 2B, 4B, and 8B. The detailed architectural hyperparameters for these configurations can be found in Appendix~Table~\ref{tab:arch}. For query initialization in the local encoder’s first cross-attention block, we implement max-pooling. Across all {\textsc BLT} models, we employ $500,000$ hashes using a single hash function, with n-gram lengths chosen between 3 and 8. All models utilize a consistent learning rate of $4e-4$. This uniformity in learning rate between token-based and {\textsc BLT} architectures is based on a hyperparameter search conducted at the 400M and 1B scales, scanning $1e-3$ to $1e-4$, which revealed $4e-4$ as optimal. For Llama-2 scaling investigations, we adopt training batch-sizes either directly as outlined in~\cite{dubey2024llama} or matched by equivalent byte count. Training employs the AdamW optimizer~\citep{adamw} with hyperparameters $\beta_1=0.9$, $\beta_2=0.95$, and $\epsilon=10^{-8}$. We warm up the learning rate linearly over 2000 iterations, then decrease it to zero using a cosine schedule. We impose a weight decay of 0.1 and apply global gradient clipping with a maximum norm of 1.0.

\section{Scaling Trends}
\label{section:scaling}

We offer a comprehensive analysis of the scaling behavior exhibited by byte-level models, providing insights that can guide the continued scaling of {\textsc BLT} models. Our scaling investigation seeks to overcome the shortcomings observed in earlier work on byte-level modeling through the following approaches: (a) We systematically examine the trends associated with compute-optimal training protocols; (b) We train 8B-parameter models using substantial training data—up to 1T tokens or 4T bytes—and assess performance on various downstream benchmarks; and (c) We analyze scaling patterns while holding inference cost constant. In a subsequent section, we delve into the distinct benefits that arise when modeling sequences at the byte level.

\subsection{Parameter Matched Compute Optimal Scaling Trends}
\label{section:parameter-matched-scaling}
We utilize the Llama 2 dataset to train several \textit{compute-optimal} bpe and {\textsc BLT} models at four distinct scales, ranging from 1B to 8B parameters. Subsequently, we graph the number of training flops versus language modeling accuracy on a representative portion of the training data mixture. The bpe models are constructed following the parameter-to-data ratio found to be optimal in Llama~3~\citep{dubey2024llama}. This approach, grounded in \textit{compute-optimal} principles, is intended to yield the highest performance possible on the training set for a specified training compute budget~\citep{hoffmann2022training}, thus serving as a solid baseline for our comparisons. Corresponding to each bpe model, we also develop a {\textsc BLT} model trained on identical data, equipping it with a Latent Transformer that is architecturally and parametrically equivalent to its bpe Transformer counterpart.

As shown in~\autoref{fig:scaling-compute-opt} (right), {\textsc BLT} models consistently perform on par with, or even better than, their bpe equivalents, and this pattern persists across increasing model capacities and computational budgets. To our knowledge, {\textsc BLT} represents the first byte-level Transformer framework to demonstrate scaling behaviors comparable to BPE-based architectures when operating in compute-efficient settings. This finding supports our hypothesis that the parameter-to-training-compute ratio identified as optimal for bpe extends to {\textsc BLT} as well, or at the very least, does not deviate significantly.

Achieving parity with bpe scaling trends necessitates both enhancements to model architectures and the adoption of dynamic patching. As shown in~\autoref{fig:scaling-compute-opt} (left), we evaluate models employing space-patching against Llama 3, approximating SpaceByte \citep{slagle2024spacebyte} by using {\textsc BLT} space-patching while omitting n-gram embeddings and cross-attention. Despite exhibiting gains over Megabyte, SpaceByte does not reach the performance levels of Llama 3. The right panel of~\autoref{fig:scaling-compute-opt} presents the cumulative benefits stemming from architectural revisions and dynamic patching. When trained at scale, {\textsc BLT} variants achieve results comparable to leading tokenizer-based models such as Llama 3.

Additionally, we find that the selection of tokenizer has a measurable impact on the efficacy of tokenizer-based approaches. Specifically, models utilizing the Llama-3 tokenizer yield better results than those employing the Llama-2 tokenizer, even when trained on identical datasets.

In summary, the {\textsc BLT} model positions itself between Llama 2 and Llama 3 in terms of performance, particularly when adopting much larger patch sizes. The bpe tokenizers used in Llama 2 and 3 correspond to average token sizes of 3.7 and 4.4 bytes, respectively. On the other hand, {\textsc BLT} achieves comparable scaling behavior with average patch sizes of 6 or even 8 bytes. Since inference flop requirements decrease as average patch size increases, selecting an 8-byte patch size can reduce inference flop by nearly 50\%. Additionally, scaling both model and data size appears to enhance the effectiveness of models with larger patch sizes. For example, with an 8-byte patch size, {\textsc BLT} initially underperforms compared to bpe Llama 2 at the 1B parameter mark but eventually surpasses the bpe variant at 7B parameters. This pattern indicates that large patch sizes might become even more advantageous as model and dataset scales continue to rise, and suggests the potential viability of employing even larger patch sizes as model capacity and available computational resources expand.

\subsection{Beyond Compute Optimal Task Evaluations}
\label{section:evals}
To further explore scaling characteristics, we train an 8B {\textsc BLT} model past the compute-optimal ratio using the {\textsc BLT}-1T dataset, which is larger and of higher quality. Model performance is then assessed across several established classification and generation benchmarks. For our evaluations, we select representative tasks focusing on common sense reasoning, world knowledge, and code generation:
\paragraph{Classification tasks} comprise ARC-Easy (0-shot)~\citep{Eval_ARC}, Arc-Challenge (0-shot)~\citep{Eval_ARC}, HellaSwag (0-shot)~\citep{Eval_hellaswag}, PIQA (0-shot)~\citep{Eval_piqa}, and MMLU (5-shot)~\citep{hendrycks2020measuring}. For these benchmarks, we use a prompt-scoring approach that computes the likelihood over the answer options, and present the mean accuracy across the tasks.
\paragraph{Coding related generation tasks:} For code generation, we report pass@1 metrics on MBPP (3-shot)~\citep{Eval_MBPP} and HumanEval (0-shot)~\citep{Eval_HumanEval}, assessing the model’s capacity to generate correct Python programs.

In \autoref{tab:evals}, we present a comparison among three models trained on the {\textsc BLT}-1T dataset: a bpe Llama 3 tokenizer-based model,\footnote{The Llama 3 tokenizer is selected due to its 128k vocabulary, which yields superior results compared to Llama 2's 32k vocabulary.} and two variations of the {\textsc BLT} model. The first uses a space-patching approach ({\textsc BLT}-Space) while the second adopts an entropy-driven patching approach ({\textsc BLT}-Entropy), incorporating an approximate monotonicity condition and resetting the context within the entropy model at every newline, as elaborated in~\autoref{section:entropy-context}. The training for all three models is standardized with the same computational budget. Additionally, for {\textsc BLT}-Entropy, we alter the entropy threshold at inference from 0.6 to 0.1, which enhances task performance though it requires additional inference steps.

The {\textsc BLT}-Entropy model achieves superior results compared to the Llama 3 model on 4 of the 7 evaluated tasks, despite being trained on an equal volume of data. This advantage likely stems from two key factors: (1) more effective utilization of training compute facilitated by dynamic patching, and (2) the explicit modeling of byte-level data rather than token-based representations.

Conversely, {\textsc BLT}-Space performs worse than the Llama 3 tokenizer on all tasks except one, though it delivers a marked decrease in inference flops owing to its higher mean patch size of 6 bytes. In contrast, both the bpe model and the entropy-patching approaches yield an average patch size of about 4.5 bytes across the training data mix. Given an identical training budget, the larger patch size enables this model to process 30\% more data than the other two approaches, which may consequently lead {\textsc BLT} to deviate further from the compute-optimal threshold.

\subsection{Patches Scale Better Than Tokens} Using {\textsc BLT} models allows for both the expansion of model parameters and the increase of patch size, all while preserving a consistent computational cost for both training and inference and holding the training data volume unchanged. The ability to scale patch size arbitrarily distinguishes patch-based architectures, permitting them to surpass the efficiency constraints typical of models relying on fixed-vocabulary tokens, as outlined in Section~\ref{section:bpe-patch}. Utilizing larger patches leads to computational savings, which can then be redirected to augment the capacity of the global latent transformer, since it requires fewer executions.

We perform a fixed inference scaling analysis to evaluate whether increasing model size, while decreasing the number of inference steps on larger patches, yields superior performance compared to employing smaller models with more steps. Using initial models with 400 million and 3.6 billion parameters and the Llama 2 tokenizer as baselines, we identify flop-matched models utilizing the Llama 3 tokenizer and {\textsc BLT}-Entropy architectures, configured with average patch sizes of 6 and 8 bytes on our training datamix (refer to~\autoref{tab:fixed-inf-params} for a summary of model configurations). For those models operating with patch size 8, three encoder layers are employed, as opposed to one in other settings. Each variant is then trained across a range of training flop budgets.

\autoref{fig:fixed_inference_scaling} demonstrates that {\textsc BLT} architectures exhibit superior scaling properties compared to tokenization-based models across both categories of inference computational costs. While BPE models initially outperform at lower training compute levels, their advantage is quickly overtaken by {\textsc BLT} models shortly after moving past the compute-optimal point. In practical scenarios, allocating additional resources to the initial pretraining phase can be advantageous, resulting in models that yield stronger performance under the same inference compute constraints. A notable illustration is found in the 8B model category, such as Llama 3.1, which was pretrained on a data volume nearly two orders of magnitude greater than the compute-optimal threshold for its model size.

The threshold at which {\textsc BLT} surpasses token-based models in performance has shifted marginally closer to the compute-optimal regime as model size increases within higher flop classes, changing from approximately 3x to 2.5x the compute-optimal budget. Additionally, the model employing the larger patch size of 8 exhibits a more pronounced scaling trajectory within the larger flop class, allowing it to outperform other configurations earlier. As elaborated in Section~\ref{section:parameter-matched-scaling}, increasing the patch size tends to result in performance that is more comparable to BPE models when scaling up model size. We hypothesize that this outcome is partly due to the relatively slower increase in the proportion of total flops allocated to the byte-level Encoder and Decoder modules compared to the Latent Transformer as the model grows. Notably, scaling the number of parameters twentyfold from 400M to 8B results in only an approximate doubling of the local model parameters for {\textsc BLT}}. This distinction is crucial because larger patch sizes only influence the flops consumed by the patch Latent Transformer, leaving the byte-level components largely unaffected. Consequently, this explains why, as model size was increased, {\textsc BLT}-Entropy with ps=8 saw its relative size shift from 1.6x to 1.7x that of the Llama 2 model.

In conclusion, our investigation into patch-length scaling reveals that the {\textsc BLT} architecture, which utilizes patches, exhibits superior scaling behavior when both the patch size and overall model size are increased together. Notably, these positive scaling dynamics continue and may even be enhanced as the model size grows.

\section{Byte Modeling Improves Robustness} We further evaluate the robustness of {\textsc BLT} against token-based counterparts that do not incorporate explicit byte-level signals, and introduce a method for converting pretrained token-based models to operate on bytes.
\label{sec:robustness}

\subsection{Character-Level Tasks}

An initial impetus for developing byte-level models was their potential to handle input noise at the byte level effectively, as well as their inherent understanding of the components within tokens—a capability where traditional tokenizer-based approaches often fall short. To quantitatively assess these aspects, we conduct supplementary experiments using benchmarks designed to test resilience to noisy input and the comprehension of character-level details across both English and multilingual contexts, encompassing not just alphabetic characters but also digits and phonemes. The findings from these evaluations are summarized in Table~\ref{tab:char_tasks}.

\paragraph{Noisy Data} To assess the robustness of tokenizer-based models in contrast to {\textsc BLT}, we generate altered versions of the benchmark classification datasets described in Section~\ref{section:evals}. For this purpose, we apply five unique character-level perturbation techniques: (a) \textit{AntSpeak}: Transforms text into a sequence of uppercase, space-separated characters. (b) \textit{Drop}: Randomly eliminates 10\% of characters from the input. (c) \textit{RandomCase}: Randomly assigns uppercase to half the characters and lowercase to the other half across the entire text. (d) \textit{Repeat}: Selects 20\% of the characters to repeat, each up to four consecutive times. (e) \textit{UpperCase}: Changes all letters to uppercase. When conducting evaluations, each of these noise schemes is individually applied to the prompt, to the completion, or simultaneously to both, thereby establishing distinct test scenarios; results are averaged over these cases. Table \ref{tab:char_tasks} summarizes performance on the noisy HellaSwag dataset~\citep{Eval_hellaswag}, revealing that {\textsc BLT} consistently surpasses tokenizer-based baselines in robustness—achieving, on average, an 8-point lead over a model trained on the same corpus, and even outperforming the Llama 3.1 model that leverages a substantially larger training set.

\paragraph{Phonology - Grapheme-to-Phoneme (G2P)} We evaluate {\textsc BLT}'s proficiency in converting sequences of graphemes—that is, the letters that make up a word—into their corresponding phonemic representations. The performance outcomes for the G2P task in a 5-shot scenario, utilizing the Phonology Bench~\citep{suvarna2024phonologybench}, are shown in Table \ref{tab:char_tasks}. Our findings demonstrate that {\textsc BLT} exceeds the results of the Llama 3 1T tokenizer-based baseline model in this setting.

\paragraph{CUTE}
To evaluate character-level reasoning, we tested {\textsc BLT} on the CUTE benchmark~\citep{edman2024cute}, which encompasses a range of tasks grouped into three main categories: compositional understanding, orthographic similarity, and sequence manipulation. CUTE represents a formidable test for conventional tokenizer-based architectures, which often demonstrate awareness of the internal makeup of their tokens, but find it challenging to leverage this insight for active text manipulation. As displayed in Table~\ref{tab:char_tasks}, {\textsc BLT}-Entropy surpasses both BPE Llama 3 variants by over 25 points on the CUTE suite. Notably, {\textsc BLT} shows outstanding effectiveness in manipulating characters, achieving an accuracy of 99.9\% on both spelling-specific tasks. These substantial gains, achieved even though {\textsc BLT} was exposed to 16 times less data than Llama 3.1, suggest that character-level patterns are particularly difficult for BPE-based models to internalize. Figure \ref{fig:cute_examples} showcases examples highlighting cases where the Llama 3 tokenizer-based system underperforms but the {\textsc BLT} model excels. Tasks involving word deletion and insertion remain the areas where BPE-based approaches retain an edge; such manipulations may not be as intuitive for byte-level frameworks. However, the performance disparity in these cases is relatively narrow, and constructing understanding from characters up to words may, in fact, prove more tractable than decomposing words back to characters. Our evaluations employed an identical protocol across all benchmarks, making use of the original prompts provided by Huggingface. It should be noted that BPE models might be improved through targeted prompt engineering.

\paragraph{Low Resource Machine Translation} We assess the performance of {\textsc BLT} on translation tasks involving twenty-one lower resource languages, which represent six major language families and employ diverse writing systems, utilizing the FLORES-101 benchmark~\citep{goyal2022flores}. SentencePiece BLEU scores are provided in Table~\ref{tab:flores}. Our evaluation shows that {\textsc BLT} consistently surpasses a counterpart trained with the Llama 3 tokenizer, providing an overall improvement of 2 BLEU points for translations into English, and a 0.5-point gain for translations out of English. For widely used language pairs, {\textsc BLT}'s performance is on par with or marginally better than that of Llama 3. Notably, for many pairs belonging to low-resource language families, {\textsc BLT} achieves superior results compared to Llama 3, highlighting the strength of byte-level modeling in handling rare and complex byte patterns.

\subsection{Training {\textsc BLT} with Llama 3 Initialization}
\label{sec:distillation}
We investigate a methodology in which {\textsc BLT} models benefit from the parameters of pre-existing, pre-trained models that utilize tokenization, thereby promoting more rapid and stable training convergence. Specifically, this is accomplished by initializing the global transformer parameters of {\textsc BLT} with those obtained from a pre-trained Llama 3.1 model. The training regime then involves updating the global transformer's parameters at a learning rate one-tenth that of the local encoder and decoder components, maintaining this schedule for the same number of steps as is optimal for Llama 3. Performance metrics are provided in Table~\ref{tab:distill}, contrasting this strategy against a standard {\textsc BLT} baseline. The results demonstrate that the {\textsc BLT} model initialized from Llama 3.1 markedly surpasses the original Llama 3 as well as the standard {\textsc BLT}, under equivalent computational constraints (in flops). Furthermore, even in comparison to our {\textsc BLT}-Entropy variant (see Table~\ref{tab:evals}), which was trained on a much larger corpus (1T tokens), the {\textsc BLT} adapted from Llama 3.1 attains higher MMLU performance, indicating that this approach can substantially decrease the compute requirements for training without sacrificing—and potentially even improving—performance.

This approach can be interpreted as converting models that rely on tokenizers into ones that do not, thereby turning a pre-trained LLaMA 3.1 model into a {\textsc BLT} model. For a thorough comparison, we report results for the original LLaMA 3.1 model, trained on 15T tokens, in Table \ref{tab:distill}, and assess its performance relative to the {\textsc BLT} version derived from LLaMA 3. Although our model shows a modest decrease in accuracy on MMLU and HumanEval benchmarks, the decline is more pronounced on several other tasks. These outcomes indicate the necessity for additional research to more effectively utilize the pre-trained model, with particular emphasis on optimizing the data mixture and related hyperparameters.

\section{Ablations and Discussion}
\label{section:ablations}
This section presents ablation studies that provide insights into the architectural decisions made for {\textsc BLT}, as well as the choices regarding patching procedures and hyper-parameter settings in the context of the {\textsc BLT} 8B parameter model trained on the {\textsc BLT}-1T dataset.

\paragraph{Entropy Model Hyper-parameters} To analyze how model size and context window length in the entropy model impact scaling behavior, we train several byte-level entropy transformer models, varying their parameter counts from 1 million up to 100 million, and testing context window sizes ranging from 64 to 512. We visualize the relationship between bits-per-byte (bpb) and training FLOPs in scaling law plots, utilizing both our $400m$ and $1b$ {\textsc BLT} models that were trained on the Llama-2 dataset; these results are shown in \autoref{fig:entropy_ablation}. Our findings indicate a positive correlation between scaling performance and both model size and context window, though gains taper off once model sizes exceed 50 million parameters.

\paragraph{Types of Patching}

We conduct an ablation study on the four patching methods described in Section~\ref{section:patching}: (1) Strided Patching utilizing strides of 4 and 6, (2) Patching determined by whitespace, (3) Patching via BPE Tokenizer relying on the Llama 3 tokenizer, and (4) Entropy-driven patching informed by a compact byte-level language model.

Although dynamic patching shortens the effective sequence length, we regulate sequence length so that all patching approaches operate over comparable context lengths. Throughout both training and inference, each model observes an equal expected number of bytes per sequence to eliminate any confounding variables related to context size differences. The outcomes of these ablation studies are presented in Figure~\ref{fig:scaling-compute-opt}. All alternative patching strategies demonstrate superior performance compared to static patching, with space patching performing nearly as well as dynamic entropy based patching.

\autoref{tab:evals_ablation} provides benchmark assessments of {\textsc BLT} models, contrasting the performances of tokenizer-based architectures, space patching, and entropy-based patching. All models were trained on the Llama 2 dataset for an optimal number of training steps as outlined by~\citet{dubey2024llama}. While space patching offers a more straightforward approach—eliminating the need to deploy an entropy model during training—we observe that the advantages demonstrated by entropy-based patching in scaling analyses (see Section~\ref{section:scaling}) extend to downstream benchmark tasks as well.\footnote{Space patching results stem from prior experiments without cross-attention, but comparable patterns appear when cross-attention is included.}

\paragraph{Cross-Attention} 
In~\autoref{tab:crossattn_ablations_1b}, we investigate the impact of incorporating cross-attention at different locations within both the encoder and decoder components of {\textsc BLT}. Regarding encoder cross-attention, we evaluate several initialization strategies for the queries: 1) employing a single learned embedding shared across all global states, 2) utilizing a hash embedding generated from the patch’s byte sequence, and 3) applying a pooling operation to the encoder’s hidden states corresponding to the patch bytes at a particular encoder layer.

Our results indicate that implementing cross-attention within the \textit{decoder} yields the greatest benefits. Incorporating cross-attention in the encoder provides modest gains, but these are limited to cases where the queries are initialized via pooling. Moreover, our analysis shows that cross-attention offers noticeable advantages on Common-Crawl data, with the effect becoming more pronounced as patch sizes increase.

\paragraph{n-gram Hash Embeddings} 

We experiment with n-gram hash embedding vocabularies of sizes 0, 100K, 200K, and 400K, and report the corresponding outcomes in Table~\ref{tab:ngram_ablations_1b}. Our observations indicate that utilizing hash embeddings confers advantages across all datasets, with the most pronounced improvements on Wikipedia and Github datasets (a difference of 0.04 bpb as opposed to 0.01 bpb after 15k steps at the 8B scale). At 8B, reducing the vocabulary from 500K to 300K hashes had a minimal impact on performance—about 0.001 bpb at 15k steps. These results suggest that hash embeddings are essential for achieving {\textsc BLT} performance comparable to that of tokenizer-based approaches. Nonetheless, we observe diminishing improvements once the number of hashes surpasses 300K. Furthermore, our findings imply that the benefits of hash embeddings are largely additive with those of cross-attention, as each contributes enhancements across distinct datasets.

\paragraph{Local Model Hyperparamaters} Table \ref{tab:local_layers} presents an ablation study exploring different configurations for the layer count in both the local encoder and decoder. Results indicate that, when used alongside hash n-gram embeddings, {\textsc BLT} achieves strong performance even with a highly simplified encoder—consisting of only a single layer—while benefiting from a more substantial decoder architecture.

\section{Related Work}
\label{section:related_work}

\textbf{Character-Level RNNs:} The task of Character Language Modeling has remained a cornerstone in neural network research since its inception~\citep{sutskever2011generating,mikolov2012subword,graves2013generating}, primarily due to its inherent capability to model novel or out-of-vocabulary words without the need for traditional back-off techniques. In their work, \cite{kim2016character} developed a model that exclusively encodes input characters using a combination of convolutional and highway networks, which subsequently provide input to LSTM-based RNNs. This architecture not only achieved results on par with then-leading RNN-based language models for English, but also demonstrated superior performance on morphologically complex languages, highlighting another important advantage of character-level large language models. Additionally, \cite{kenter2018byte} introduced byte-level LSTM models for machine comprehension tasks and observed that these outperformed their word-level counterparts, particularly in morphologically rich languages such as Turkish and Russian. Similarly, \cite{zhang2015character} investigated classification using character-based convolutional models and found these models to surpass word-level models in several tasks. Further advances were made by \citet{chung2019hierarchical}, who employed hierarchical LSTM models equipped with boundary detectors at each hierarchical stage, uncovering latent text structure and thus further optimizing character-level language modeling. Distinct from attention-based frameworks, ByteNet~\cite{kalchbrenner2016neural} utilizes convolutional neural network layers applied directly at the character level for tasks such as machine translation.

\textbf{Character-Level Transformers:} Advances in transformer architectures leveraging attention mechanisms~\citep{vaswani2017attention}, in tandem with the adoption of subword tokenization~\citep{sennrich-etal-2016-neural}, have greatly enhanced neural approaches to language modeling and evaluations on benchmark datasets. Nonetheless, using word and subword tokens inherently prescribes the abstraction layer at which these models function. Aiming to synthesize the advances offered by transformer networks with early favorable outcomes in character-level language modeling, \citet{al2019character} explored training very deep transformer models supplemented by auxiliary objectives. These transformer-based systems surpassed previous character language models built on LSTM architectures, yet a considerable performance disparity with word-level LLMs remained. Additionally, results from GPT-2~\citep{radfordgpt2} indicated that on extensive corpora such as the 1 Billion Word Benchmark, byte-level language models did not reach parity with those built on word-level inputs.

\cite{Choe2019BridgingTG} showed that transformer-based language models trained at the byte level can surpass subword-level LLMs with similar numbers of parameters in terms of performance. However, these byte-level models require significantly more computational resources and extended training periods. In a related effort, \cite{el2020characterbert} introduced CharFormer, a variant of BERT that generates word-level representations through convolutional operations over character embeddings. Although CharFormer demonstrated gains in medical-domain benchmarks, these results also came at the cost of substantially increased compute usage. \cite{clark2022canine} proposed CANINE, an encoder-only model with 150M parameters that operates directly on character sequences. This architecture integrates a deep transformer backbone—akin to the global module in our framework—and combines it with both local transformers and strided convolutions to reduce the character sequence length. CANINE achieves superior performance to mBERT, a token-based encoder-only model, on various multilingual tasks. In the realm of encoder-decoder architectures, ByT5~\citep{xue2022byt5} explored the viability of byte-level modeling without employing input patching. While ByT5 proved more resilient to noise and rivaled tokenized models using just a quarter of the dataset, the absence of input reduction mechanisms meant it was forced to compute attention across every input byte, leading to substantial compute inefficiency. The inherently longer input sequences that accompany byte-level modeling present formidable obstacles for scaling such models efficiently. More recently, with the advent of the Mamba Architecture~\citep{gu2023mamba}, which sustains a fixed-sized memory footprint even for very long contexts, \cite{wang2024mambabyte} developed a byte-level Mamba model that forgoes patching and was shown, in compute-matched experiments at the 350M parameter scale, to surpass byte-level transformers in bits-per-byte across several datasets.

\textbf{Patching-based approaches:} Utilizing patching techniques has the potential to significantly reduce the excessive computational demands typically associated with byte-level LLMs, while potentially preserving their effectiveness. Several studies have shown promising results at smaller scales in terms of both model size and the volume of training data. For instance, \cite{nawrot-etal-2022-hierarchical} introduce the hourglass transformer, which incorporates static patching through downsampling and upsampling; their approach demonstrates superior performance compared to other byte-level methods at the 150M parameter scale. Building on this, \cite{nawrot-etal-2023-efficient} propose enhancements using dynamic patching methods. These include an end-to-end trainable boundary predictor, a version of the predictor supervised with specific tokenizers, and an entropy-informed patching approach akin to {\textsc BLT}. They report that, on a dataset of approximately 400M tokens with 40M-parameter models, these mechanisms can surpass the language modeling performance of standard transformers of that era. Furthermore, \cite{lester2024training} explore training on input sequences that are compressed through arithmetic coding, achieving higher compression rates than those realized by BPE. By employing an equal-info windows strategy, their method surpasses byte-level baselines on language modeling tasks, though it does not outperform baselines based on subword tokenization.

This study is primarily inspired by and most directly connected to the MegaByte model~\citep{yu2023megabyte}, a causal decoder-only large language model that employs a static, pre-defined patching and concatenation scheme to map bytes to patches, while its decoder utilizes a local model to reconstruct bytes from these patches. MegaByte has been shown to achieve comparable performance to tokenizer-based approaches at the 1B parameter level on datasets amounting to approximately 400B bytes. In our experiments, we systematically evaluate MegaByte and observe that its static patching framework falls short of the current state-of-the-art, compute-optimal tokenizer-based models under fixed computation constraints. Furthermore, we illustrate how {\textsc BLT} effectively narrows this performance discrepancy. The same limitation is highlighted by \cite{slagle2024spacebyte}, who propose enhancing static patching to account for whitespaces and other space-character bytes, and incorporate an additional local encoder model. Their results indicate certain gains relative to tokenizer-based Transformers within a compute-constrained environment, particularly on specialized datasets like Github and arXiv with models at the 1B parameter scale. We further include experimental analysis of this approach and reveal that more substantial architectural advances will be necessary to enable byte-level models to reach parity with current leading token-based architectures, including models such as Llama 3.

\section{Limitations and Future Work}

For the selection of architectural parameters in this study, we train models for the empirically determined optimal number of steps utilized in Llama~3~\citep{dubey2024llama}. It should be noted, however, that these scaling principles were derived with BPE-level transformers in mind, which may yield suboptimal ratios of data and parameter sizes for {\textsc BLT}. We defer the formulation of scaling laws specifically tailored to {\textsc BLT}—a development that could potentially reveal even more advantageous scaling behavior for our model architecture—to future investigation. Furthermore, a significant portion of our experimental evaluation has been performed at parameter counts up to 1B; as we scale models to 8B parameters or larger, the ideal architectural configurations may shift, which could enable greater improvements in large-scale performance.

Current transformer codebases and libraries are predominantly optimized for transformer architectures that utilize tokenizers. Although our study includes theoretically equivalent flop experiments and integrates select efficient solutions like FlexAttention for non-standard transformer components, our actual implementations may still lag behind tokenizer-based models regarding real-world runtime performance, suggesting that there remains room for additional optimization.

Whereas {\textsc BLT} employs a distinct, pre-trained entropy model for patching, integrating the training of the patching model into a fully end-to-end framework presents a promising avenue for future exploration. Section \ref{sec:distillation} details preliminary experiments that offer encouraging results for converting tokenizer-based models, like Llama 3—which is pretrained on over 10T tokens—into byte-based models, by transferring and freezing the global transformer weights. Continued investigation could yield strategies that maintain the advantages of bytefying, while potentially surpassing the capabilities of the original tokenizer-based architectures, all without necessitating training from the beginning.

\section{Conclusion}
\label{section:conclusion}

In this work, we introduce the Byte Latent Transformer (\textbf{BLT}), an innovative model architecture that moves away from the standard reliance on fixed-vocabulary tokenization in large language models. By employing an adaptive, trainable mechanism that segments bytes into patches, {\textsc BLT} enables a more efficient allocation of computation aligned with the underlying complexity of the input data. This approach yields notable gains in both computational efficiency and model robustness. Our comprehensive scaling experiments reveal that {\textsc BLT} can achieve comparable results to established tokenization-based models such as Llama 3, at scales up to 8B parameters and 4T bytes of data. Moreover, {\textsc BLT} facilitates up to a 50\% reduction in inference flops with only marginal impact on common evaluation metrics. Importantly, {\textsc BLT} introduces a novel scaling paradigm, supporting parallel increases in both model capacity and patch size without exceeding fixed inference resource constraints—a substantial benefit for practical deployment environments. Beyond processing raw byte-level inputs, {\textsc BLT} enhances the model's capacity to manage infrequent data phenomena, yielding greater robustness to noisy inputs and improved representation of sub-word granularity. Altogether, these advances position {\textsc BLT} as a compelling, scalable alternative to traditional tokenization-driven methods, offering a robust and efficient basis for next-generation language model development.

\end{document}